% ============================================================================
% RAPPORT ACADÉMIQUE — Système de Prise de Rendez-vous Médicaux
% Compilable avec : pdflatex rapport.tex
% ============================================================================

\documentclass[12pt, a4paper]{article}

% --- Packages ---------------------------------------------------------------
\usepackage[utf8]{inputenc}          % Encodage UTF-8
\usepackage[T1]{fontenc}             % Encodage des polices
\usepackage[french]{babel}           % Typographie française
\usepackage{geometry}                % Mise en page
\usepackage{graphicx}                % Images
\usepackage{hyperref}                % Liens cliquables
\usepackage{xcolor}                  % Couleurs
\usepackage{array}                   % Tableaux améliorés
\usepackage{booktabs}                % Lignes horizontales dans les tableaux
\usepackage{fancyhdr}                % En-têtes et pieds de page
\usepackage{titlesec}                % Personnalisation des titres
\usepackage{setspace}                % Interligne
\usepackage{parskip}                 % Espacement entre paragraphes
\usepackage{enumitem}                % Listes personnalisées

% --- Géométrie de la page ---------------------------------------------------
\geometry{
  top=2.5cm,
  bottom=2.5cm,
  left=2.5cm,
  right=2.5cm
}

% --- Couleurs personnalisées -------------------------------------------------
\definecolor{fssm-blue}{RGB}{0, 74, 153}
\definecolor{fssm-dark}{RGB}{30, 40, 60}

% --- Hyperliens --------------------------------------------------------------
\hypersetup{
  colorlinks=true,
  linkcolor=fssm-blue,
  urlcolor=fssm-blue,
  citecolor=fssm-blue,
  pdfauthor={Étudiants FSSM},
  pdftitle={Système de Prise de Rendez-vous Médicaux},
  pdfsubject={Rapport de projet — Génie Logiciel},
}

% --- En-têtes / pieds de page -----------------------------------------------
\pagestyle{fancy}
\fancyhf{}
\fancyhead[L]{\small\textit{Système de Prise de Rendez-vous Médicaux}}
\fancyhead[R]{\small\textit{FSSM — Génie Logiciel}}
\fancyfoot[C]{\thepage}
\renewcommand{\headrulewidth}{0.4pt}

% --- Style des sections ------------------------------------------------------
\titleformat{\section}
  {\Large\bfseries\color{fssm-blue}}{\thesection}{1em}{}
\titleformat{\subsection}
  {\large\bfseries\color{fssm-dark}}{\thesubsection}{0.8em}{}

% --- Interligne --------------------------------------------------------------
\setstretch{1.15}

% ============================================================================
% DÉBUT DU DOCUMENT
% ============================================================================
\begin{document}

% ============================================================================
% PAGE DE GARDE
% ============================================================================
\begin{titlepage}
  \centering
  \vspace*{1cm}

  % Université
  {\Large\textbf{Université Cadi Ayyad}}\\[4pt]
  {\large Faculté des Sciences Semlalia — Marrakech}\\[4pt]
  {\large Département d'Informatique}\\[1.5cm]

  % Ligne décorative
  {\color{fssm-blue}\rule{0.7\textwidth}{1.5pt}}\\[1.5cm]

  % Titre
  {\LARGE\bfseries\color{fssm-blue} Système de Prise de Rendez-vous\\[6pt] Médicaux en Ligne}\\[1.2cm]

  {\color{fssm-blue}\rule{0.7\textwidth}{1.5pt}}\\[1.5cm]

  % Sous-titre
  {\large Rapport de projet — Module de Génie Logiciel}\\[2cm]

  % Informations
  \begin{tabular}{rl}
    \textbf{Réalisé par :}  & Étudiants FSSM \\[6pt]
    \textbf{Encadré par :}  & Pr. Encadrant \\[6pt]
    \textbf{Filière :}      & Sciences Mathématiques et Informatique \\[6pt]
    \textbf{Année :}        & 2025 / 2026 \\
  \end{tabular}

  \vfill
  {\small Marrakech — Juin 2026}
\end{titlepage}

% ============================================================================
% TABLE DES MATIÈRES
% ============================================================================
\tableofcontents
\newpage

% ============================================================================
% 1. INTRODUCTION
% ============================================================================
\section{Introduction}

\subsection{Contexte}

La gestion des rendez-vous médicaux demeure un défi majeur dans les établissements de santé au Maroc. Les méthodes traditionnelles — appels téléphoniques, files d'attente physiques — engendrent des pertes de temps considérables tant pour les patients que pour le personnel soignant. Les oublis de rendez-vous, les doublons et l'absence de traçabilité numérique des dossiers médicaux aggravent la situation.

Dans ce contexte, le recours à une plateforme web dédiée offre une réponse concrète à ces problématiques. Le présent projet s'inscrit dans le cadre du module de \textbf{Génie Logiciel} dispensé à la Faculté des Sciences Semlalia de Marrakech (FSSM) et vise à appliquer les méthodologies de conception et de développement logiciel à un cas d'usage réel.

\subsection{Objectifs du projet}

Le projet poursuit les objectifs suivants :

\begin{itemize}[nosep]
  \item Permettre aux patients de réserver, modifier et annuler des rendez-vous en ligne.
  \item Fournir aux médecins un tableau de bord complet pour gérer leur planning et consulter les dossiers patients.
  \item Mettre en place un système de notifications internes pour le suivi des rendez-vous.
  \item Intégrer un mécanisme de partage du dossier médical via QR code.
  \item Garantir la sécurité des données grâce à une authentification par jetons JWT.
\end{itemize}

% ============================================================================
% 2. ANALYSE DES BESOINS
% ============================================================================
\section{Analyse des besoins}

\subsection{Problématique}

Comment concevoir une plateforme web ergonomique et sécurisée qui centralise la prise de rendez-vous médicaux, la gestion du planning des médecins et le suivi des dossiers patients, tout en offrant une expérience utilisateur adaptée aux différents rôles (patient, médecin, administrateur) ?

\subsection{Exigences fonctionnelles}

Le tableau~\ref{tab:ef} synthétise les principales exigences fonctionnelles identifiées lors de la phase d'analyse.

\begin{table}[h!]
\centering
\caption{Exigences fonctionnelles du système}
\label{tab:ef}
\small
\begin{tabular}{>{\bfseries}p{1cm} p{4cm} p{8cm}}
\toprule
ID & Fonctionnalité & Description \\
\midrule
EF1 & Inscription / Connexion & Création de compte avec rôle (patient ou médecin), authentification sécurisée par JWT. \\
EF2 & Gestion des rendez-vous & Création, consultation, modification, annulation et suivi de l'état (en attente, confirmé, annulé, terminé). \\
EF3 & Dashboard patient & Vue des prochains rendez-vous, recherche par spécialité ou médecin, accès au dossier médical. \\
EF4 & Dashboard médecin & Planning du jour, liste des patients, confirmation des rendez-vous, accès au dossier patient par QR code. \\
EF5 & Notifications & Alertes internes pour les nouveaux rendez-vous, confirmations, annulations et rappels. \\
EF6 & Dossier médical & Informations médicales, signes vitaux, historique des consultations et documents médicaux. \\
EF7 & QR Code & Génération et téléchargement d'un QR code par le patient, scan par le médecin pour accéder au dossier. \\
\bottomrule
\end{tabular}
\end{table}

\subsection{Exigences non fonctionnelles}

\begin{table}[h!]
\centering
\caption{Exigences non fonctionnelles}
\label{tab:enf}
\small
\begin{tabular}{>{\bfseries}p{1cm} p{3.5cm} p{8.5cm}}
\toprule
ID & Critère & Description \\
\midrule
ENF1 & Sécurité & Authentification JWT, hashage des mots de passe, contrôle d'accès par rôle. \\
ENF2 & Performance & Temps de réponse de l'API inférieur à 500\,ms pour les opérations courantes. \\
ENF3 & Portabilité & Conteneurisation via Docker pour un déploiement reproductible. \\
ENF4 & Ergonomie & Interface responsive adaptée au mobile et au bureau, avec navigation intuitive. \\
ENF5 & Maintenabilité & Architecture modulaire en micro-services, code structuré et documenté. \\
\bottomrule
\end{tabular}
\end{table}

% ============================================================================
% 3. CONCEPTION
% ============================================================================
\section{Conception}

\subsection{Architecture générale}

L'application adopte une architecture \textbf{client-serveur} découpée en trois services principaux, communiquant via des API REST :

\begin{table}[h!]
\centering
\caption{Services de l'architecture}
\label{tab:archi}
\small
\begin{tabular}{>{\bfseries}p{4cm} p{9cm}}
\toprule
Service & Responsabilité \\
\midrule
Auth Service (\texttt{users}) & Inscription, connexion, gestion des profils (patient, médecin), émission et rafraîchissement des jetons JWT. \\
Appointment Service & Opérations CRUD sur les rendez-vous, statistiques du tableau de bord, filtrage par date et statut. \\
Notification Service & Création, listing et marquage en lu des notifications liées aux rendez-vous. \\
\bottomrule
\end{tabular}
\end{table}

Le \textbf{frontend} React communique avec le backend Django REST via un client HTTP Axios configuré avec un intercepteur d'authentification. Les tokens JWT sont stockés dans le \texttt{localStorage} et automatiquement injectés dans les en-têtes des requêtes.

\subsection{Diagrammes UML}

Les diagrammes suivants ont été élaborés pour modéliser le système :

\textbf{Diagramme de cas d'utilisation.} Il identifie trois acteurs principaux — \textit{Patient}, \textit{Médecin} et \textit{Administrateur} — et leurs interactions avec le système : prise de rendez-vous, gestion du planning, consultation du dossier médical et administration des utilisateurs.

\textbf{Diagramme de classes.} Le modèle de données s'articule autour de cinq entités principales : \texttt{User} (avec héritage d'\texttt{AbstractUser}), \texttt{DoctorProfile}, \texttt{PatientProfile}, \texttt{Appointment} et \texttt{Notification}. Deux entités complémentaires enrichissent le dossier médical : \texttt{VitalSigns} et \texttt{MedicalDocument}. Les relations clés sont les suivantes :

\begin{itemize}[nosep]
  \item \texttt{User} $\longleftrightarrow$ \texttt{DoctorProfile} / \texttt{PatientProfile} : relation \textit{OneToOne}.
  \item \texttt{User} $\longleftrightarrow$ \texttt{Appointment} : double clé étrangère (\textit{patient}, \textit{doctor}).
  \item \texttt{User} $\longleftrightarrow$ \texttt{Notification} : relation \textit{ManyToOne}.
  \item \texttt{User} $\longleftrightarrow$ \texttt{VitalSigns} / \texttt{MedicalDocument} : relation \textit{ManyToOne}.
\end{itemize}

\textbf{Diagramme de séquence — Prise de rendez-vous.} Le patient recherche un médecin, sélectionne un créneau, soumet le formulaire. Le backend valide la requête, crée l'objet \texttt{Appointment} avec le statut \texttt{pending}, puis génère une notification pour le médecin concerné.

% ============================================================================
% 4. IMPLÉMENTATION
% ============================================================================
\section{Implémentation}

\subsection{Technologies utilisées}

\begin{table}[h!]
\centering
\caption{Stack technologique}
\label{tab:tech}
\small
\begin{tabular}{>{\bfseries}p{3cm} p{3.5cm} p{6.5cm}}
\toprule
Couche & Technologie & Rôle \\
\midrule
Backend & Django 6.0 & Framework web Python pour la logique métier. \\
API REST & DRF 3.17 & Sérialisation, routage et permissions des endpoints. \\
Authentification & SimpleJWT 5.5 & Émission et vérification des jetons JWT. \\
Base de données & SQLite (dev) & Stockage relationnel léger pour le développement. \\
Frontend & React 18 & Construction de l'interface utilisateur (SPA). \\
Routage client & React Router 6 & Navigation côté client avec routes protégées. \\
HTTP Client & Axios & Requêtes HTTP avec intercepteur JWT. \\
QR Code & qrcode.react & Génération de QR codes côté client. \\
Conteneurisation & Docker Compose & Orchestration des services backend, frontend et Redis. \\
\bottomrule
\end{tabular}
\end{table}

\subsection{Description des principales fonctionnalités}

\textbf{Authentification et autorisation.} Le module \texttt{users} implémente l'inscription avec attribution automatique du rôle et la connexion avec génération de paires de jetons (access / refresh). Un composant \texttt{PrivateRoute} côté React protège les routes sensibles et redirige l'utilisateur selon son rôle vers le tableau de bord approprié.

\textbf{Gestion des rendez-vous.} Le modèle \texttt{Appointment} gère un cycle de vie à quatre états : \textit{pending} $\rightarrow$ \textit{confirmed} $\rightarrow$ \textit{completed} ou \textit{cancelled}. Le sérialiseur valide que le médecin sélectionné possède bien le rôle \texttt{doctor}. Les vues filtrent automatiquement les rendez-vous selon le rôle de l'utilisateur authentifié.

\textbf{Tableaux de bord différenciés.} Le composant \texttt{DashboardPatient} propose trois onglets — accueil avec recherche, agenda et profil — tandis que \texttt{DashboardDoctor} affiche le planning du jour, la liste des patients et les statistiques. La sélection du tableau de bord est effectuée dynamiquement par le composant \texttt{DashboardRouter} en fonction du rôle.

\textbf{Dossier médical et QR code.} Le patient peut consulter son dossier complet (informations médicales, signes vitaux, historique, documents) et générer un QR code contenant son identifiant. Le médecin scanne ce QR code pour accéder au dossier via l'endpoint \texttt{/users/patients/<id>/}, protégé par une vérification de rôle.

\textbf{Notifications.} Le service de notifications crée des alertes pour chaque événement lié aux rendez-vous (nouveau, confirmé, annulé, rappel). Les notifications sont récupérées et marquées comme lues via des endpoints dédiés.

% ============================================================================
% 5. TESTS ET VALIDATION
% ============================================================================
\section{Tests et validation}

\subsection{Méthodes de test}

La stratégie de test adoptée couvre trois niveaux complémentaires :

\begin{table}[h!]
\centering
\caption{Stratégie de test}
\label{tab:tests}
\small
\begin{tabular}{>{\bfseries}p{3.5cm} p{3cm} p{6.5cm}}
\toprule
Type de test & Outil & Périmètre \\
\midrule
Tests unitaires (backend) & Django TestCase & Modèles, sérialiseurs, vues (users, appointments, notifications). \\
Tests d'API & DRF APIClient & Endpoints REST : codes HTTP, validation des données, contrôle d'accès. \\
Tests frontend & React Testing Library & Rendu des composants, interactions utilisateur de base. \\
\bottomrule
\end{tabular}
\end{table}

\subsection{Résultats obtenus}

La suite de tests backend comprend plus de \textbf{25 cas de test} répartis sur trois fichiers, couvrant les scénarios suivants :

\begin{itemize}[nosep]
  \item \textbf{Rendez-vous} : création avec médecin valide/invalide, listing filtré par rôle, modification du statut, suppression, isolation des données entre patients, rendez-vous du jour et statistiques du tableau de bord.
  \item \textbf{Utilisateurs} : inscription (patient et médecin), connexion avec identifiants valides/invalides, accès au profil, listing des médecins avec filtrage par spécialité.
  \item \textbf{Notifications} : listing des notifications de l'utilisateur connecté, marquage en lu, isolation entre utilisateurs.
\end{itemize}

L'ensemble des tests passent avec succès, confirmant la conformité du système aux exigences fonctionnelles spécifiées.

% ============================================================================
% 6. CONCLUSION ET PERSPECTIVES
% ============================================================================
\section{Conclusion et perspectives}

Ce projet a permis de concevoir et d'implémenter une plateforme complète de prise de rendez-vous médicaux, couvrant l'ensemble du cycle de vie : de l'inscription des utilisateurs à la gestion des consultations, en passant par le suivi du dossier médical et le système de notifications. L'architecture modulaire adoptée — avec Django REST Framework en backend et React en frontend — garantit une séparation claire des responsabilités et facilite l'évolution du système.

L'intégration du partage de dossier médical par QR code constitue une valeur ajoutée notable, offrant un accès rapide et sécurisé aux informations médicales lors des consultations. La conteneurisation via Docker assure la portabilité et la reproductibilité du déploiement.

Plusieurs perspectives d'amélioration se dessinent pour les versions futures :

\begin{itemize}[nosep]
  \item \textbf{Notifications par e-mail et SMS} pour les rappels de rendez-vous, en exploitant le service Redis déjà intégré.
  \item \textbf{Système de paiement en ligne} pour le règlement des consultations.
  \item \textbf{Téléconsultation} via intégration d'un module de visioconférence (WebRTC).
  \item \textbf{Migration vers PostgreSQL} pour un environnement de production robuste.
  \item \textbf{Application mobile native} (React Native) pour améliorer l'accessibilité.
  \item \textbf{Tableau de bord administrateur} avancé avec analyses statistiques et gestion globale des utilisateurs.
\end{itemize}

% ============================================================================
% FIN DU DOCUMENT
% ============================================================================
\end{document}
